%% toc: P15 172 172.3,172.4
% Huber 1964: general-score contamination bias, the asymmetric variance correction, the standing of the unknown-scale proposals, and the pair-average median under arbitrary symmetric contamination.
\chapter{Huber's General Bias Bound, Proposal Standing and Pair-Median Risk}
\label{sec:p15-w02b-p15-071-huber-completions}

This chapter completes three parts of Huber's 1964 discussion that
\cref{sec:p15-huber-remaining-methods} records only for special cases.
The first concerns the bias of a bounded monotone score when the
contamination is not symmetric: Huber derives a mean-value identity
and a bound for an arbitrary such score, and asserts that the extra
variance caused by asymmetric contaminants is negligible beside the bias
\cite[Sec.~7, pp.~82--83]{Huber1964RobustLocation}.
The second concerns the asymptotic claims attached to his three
unknown-scale proposals \cite[Sec.~11, pp.~96--97]{Huber1964RobustLocation}.
The third extends his bound for the pair-average median from contaminants
with square-integrable densities to every symmetric contaminant
\cite[Sec.~12, pp.~100--101]{Huber1964RobustLocation}.
The inputs are a symmetric central law, a contamination fraction and
a bounded score; the conclusions are bias bounds, worst-case limiting
variances and their attainment. Results for the general information game
are in \cref{sec:p15-huber-information-game}.

\section{The general contamination-bias identity}
\label{sec:p15-w02b-p15-071-bias}

\begin{hbdefinition}{def:p15-w02b-p15-071-bias-setting}{Asymmetric contamination}
Let \(G\) be a probability on \(\mathbb R\) symmetric about zero,
let \(H\) be any probability on \(\mathbb R\), let
\(0\leq\varepsilon<1\), and put \(F=(1-\varepsilon)G+\varepsilon H\).
Let \(\psi:\mathbb R\to\mathbb R\) be bounded, nondecreasing and
skew symmetric, \(\psi(-x)=-\psi(x)\), and write
\(s=\sup_x|\psi(x)|\).
For any probability \(K\) on \(\mathbb R\), put
\[
 \lambda_K(t)=E_K\psi(X-t),
\]
which is finite and nonincreasing in \(t\).
A real \(c=c(F)\) with \(\lambda_F(c)=0\) is a population target of
the \(\psi\)-estimate under \(F\); Huber calls it the bias caused by
the contamination \(\varepsilon H\). Existence of such a zero is a
premise. If the zeros form an interval, each statement below applies
to each zero.
Huber's suggested reading is that most observations follow
\(G(x-\vartheta)\) while a fraction \(\varepsilon\) follows a law
unrelated to the location \(\vartheta\)
\cite[p.~82]{Huber1964RobustLocation}.
\end{hbdefinition}

\begin{hbtheorem}{thm:p15-w02b-p15-071-bias-identity}{proved=annex:p15-w02b-p15-071-bias-identity}{The mean-value bias identity and bound (2)}
In the setting of \cref{def:p15-w02b-p15-071-bias-setting}:
\begin{enumerate}
\item \(\lambda_G(0)=0\), \(\lambda_F(0)=\varepsilon\lambda_H(0)\)
 and \(|\lambda_H(0)|\leq s\).
\item Suppose \(\lambda_G\) and \(\lambda_H\) are differentiable on an
 open interval containing the closed segment between \(0\) and \(c\).
 Then there is \(\vartheta\in(0,1)\) with
 \(\lambda_F(0)+c\lambda_F'(\vartheta c)=0\).
 If \(\lambda_G'(\vartheta c)<0\), the denominator below is negative and
 \[
  c(F)=-\frac{\varepsilon\lambda_H(0)}
  {(1-\varepsilon)\lambda_G'(\vartheta c)+\varepsilon\lambda_H'(\vartheta c)},
 \qquad
  |c(F)|\leq\frac{\varepsilon}{1-\varepsilon}\,
              \frac{s}{-\lambda_G'(\vartheta c)} .
 \]
 The second inequality is the source's bound (2).
\item Suppose only \(\lambda_G\) is differentiable on such an interval.
 Then there is \(\vartheta'\in(0,1)\) with
 \((1-\varepsilon)c\lambda_G'(\vartheta'c)=-\varepsilon\lambda_H(c)\),
 and, whenever \(\lambda_G'(\vartheta'c)<0\),
 \[
  |c(F)|\leq\frac{\varepsilon}{1-\varepsilon}\,
              \frac{s}{-\lambda_G'(\vartheta'c)} .
 \]
 No derivative of \(\lambda_H\) is used; atoms of \(H\) at the
 discontinuities of \(\psi\) are allowed.
\item Suppose \(\lambda_G\) is differentiable on a neighborhood of zero,
 its derivative is continuous at zero, and \(\lambda_G'(0)<0\).
 Let \(\overline c_\varepsilon\) be the supremum of \(|c|\) over all
 probabilities \(H\) and all zeros \(c\) of the resulting \(\lambda_F\)
 (zero if there are none). Then \(\overline c_\varepsilon\to0\) as
 \(\varepsilon\downarrow0\), and
 \[
  \overline c_\varepsilon\leq
  \frac{\varepsilon}{1-\varepsilon}\,\frac{s}{-\lambda_G'(0)}\,(1+o(1)).
 \]
 This is the precise form of replacing \(\vartheta c\) by zero for
 small \(\varepsilon\).
\end{enumerate}
\end{hbtheorem}

Part (iii) is added here; the source states the identity and the bound
with the same intermediate point \cite[p.~83]{Huber1964RobustLocation}.
The bound is uniform over all contaminants, including those placing
mass where \(\psi\) jumps.

\begin{hbexample}{def:p15-w02b-p15-071-bias-normal}{The clipped score at the normal center}
Take \(G=\Phi\) and \(\psi=\psi_k\) of
\cref{def:p15-huber-location-score}. Then \(s=k\),
\(\lambda_\Phi(t)=m_k(t)\) of \cref{thm:p15-huber-asymmetric-bias},
and \(-\lambda_\Phi'(t)=\Phi(t+k)-\Phi(t-k)\), which is continuous with
value \(A(k)=2\Phi(k)-1\) at zero. The first-order bound of part (iv) is
\[
 \frac{\varepsilon k}{(1-\varepsilon)A(k)} .
\]
At \((\varepsilon,k)=(0.1,1.14)\), \((0.01,1.95)\), \((0.001,2.64)\)
it equals approximately \(0.16986\), \(0.020759\), \(0.0026647\).
These round to the source's values \(0.17\), \(0.021\), \(0.0027\)
\cite[p.~83]{Huber1964RobustLocation}. The three thresholds are the
normalized constants of \cref{thm:p15-huber-least-favorable-normal},
which are approximately \(1.14017\), \(1.94511\), \(2.63288\).
The exact maximal biases \(c_*\) solving
\(m_k(c_*)=-\varepsilon k/(1-\varepsilon)\) are approximately
\(0.17038\), \(0.0207597\), \(0.00266474\); by
\cref{thm:p15-huber-asymmetric-bias} they are attained when \(H\) puts
all its mass to the right of \(k+c_*\). The first-order values are
therefore approximations, slightly below the attained maxima.
The sample sizes \(1/(4c^2)\) at the first-order values are approximately
\(8.66\), \(580.1\), \(35207\); the source reports \(9\), \(600\)
and \(35000\).
\end{hbexample}

\begin{hbtheorem}{thm:p15-w02b-p15-071-bias-variance}{proved=annex:p15-w02b-p15-071-bias-variance}{Admitting asymmetric contaminants changes the worst variance by o(c)}
In the setting of \cref{def:p15-w02b-p15-071-bias-setting}, assume also
that \(\psi\) is continuous and not identically zero, and let the
estimate be the midpoint minimizer of \(\sum_i\rho(x_i-t)\),
\(\rho(u)=\int_0^u\psi\). Assume \(\lambda_G\) is differentiable on a
neighborhood \(N\) of zero, and put
\[
 a_G(t)=E_G\psi(X-t)^2,\qquad b_G(t)=-\lambda_G'(t)\quad(t\in N).
\]
Assume \(a_G\) and \(b_G\) are differentiable at zero and \(b_G(0)>0\).
For \(0\leq\varepsilon<1\), let \(\mathcal H_\varepsilon\) be the set of
probabilities \(H\) for which \(\lambda_F\) has a zero \(c\in N\) with
\(b_G(c)>0\) at which \(\lambda_H\) is differentiable. Put
\[
 R_\varepsilon(t)=\frac{(1-\varepsilon)a_G(t)+\varepsilon s^2}
                       {(1-\varepsilon)^2b_G(t)^2}.
\]
\begin{enumerate}
\item For \(H\in\mathcal H_\varepsilon\), the zero \(c\) is unique,
 the estimate is consistent for \(c\), and
 \[
  \sqrt n(T_n-c)\Longrightarrow N(0,V(\psi,F)),\qquad
  V(\psi,F)=\frac{E_F\psi(X-c)^2}{\lambda_F'(c)^2}\leq R_\varepsilon(c).
 \]
\item The functions \(a_G,b_G\) are even, so their derivatives vanish at
 zero. There is \(\eta(\tau)\to0\) as \(\tau\downarrow0\) such that
 \(R_\varepsilon(t)\leq R_\varepsilon(0)+\tau\eta(\tau)\) whenever
 \(0\leq\varepsilon\leq1/2\) and \(|t|\leq\tau\).
 If also \(\lambda_G'\) is continuous at zero, then for all sufficiently
 small \(\varepsilon\),
 \[
  \sup_{H\in\mathcal H_\varepsilon}V(\psi,F)\leq
  R_\varepsilon(0)+\overline c_\varepsilon\,\eta(\overline c_\varepsilon),
 \]
 with \(\overline c_\varepsilon\to0\) as in
 \cref{thm:p15-w02b-p15-071-bias-identity}. The correction is
 \(o(\overline c_\varepsilon)\).
\item For \(G=\Phi\) and \(\psi=\psi_k\), every \(H\) with
 \(H\{c-k,c+k\}=0\) at the zero \(c\) of \(\lambda_F\) belongs to
 \(\mathcal H_\varepsilon\), \(R_\varepsilon(0)=V^*(k,\varepsilon)\) of
 \cref{thm:p15-huber-normal-variance}, and for all sufficiently small
 \(\varepsilon\)
 \[
  0\leq\sup_{H\in\mathcal H_\varepsilon}V(\psi_k,F)-V^*(k,\varepsilon)
   \leq\overline c_\varepsilon\,\eta(\overline c_\varepsilon).
 \]
\end{enumerate}
\end{hbtheorem}

The theorem supplies explicit regularity conditions for the source's
statement that the correction is of order \(o(c)\)
\cite[p.~83]{Huber1964RobustLocation}. It does not cover contaminants
with an atom at \(c\pm k\) for the clipped score: there the limit is the
two-half-normal law of \cref{thm:p15-convex-location-normality}, whose
variance is not the displayed ratio. It also does not assert that the
bias and the variance correction together give a mean squared error;
the bias is a displacement of the population target and does not shrink
with \(n\).

\section{Standing of the unknown-scale proposals}
\label{sec:p15-w02b-p15-071-proposals}

Huber states that Proposals 1 and 2 of
\cref{def:p15-huber-three-proposals} are asymptotically minimax when
\(k\) is chosen in relation to the true \(\varepsilon\), and that all
three proposals are asymptotically equivalent to Winsorizing, differing
in how the number of Winsorized observations is determined
\cite[p.~96]{Huber1964RobustLocation}. The following results give each
claim its exact standing.

\begin{hbtheorem}{thm:p15-w02b-p15-071-winsor}{proved=annex:p15-w02b-p15-071-winsor}{Each proposal's location is a Winsorized mean}
Let \(S>0\), \(k>0\), and let \(T\) satisfy
\(\sum_i\psi(k,(x_i-T)/S)=0\). Replace each \(x_i<T-kS\) by \(T-kS\)
and each \(x_i>T+kS\) by \(T+kS\). Then \(T\) is the arithmetic mean of
the modified sample. This holds for every solution of Proposals 1 and 2,
and for Proposal 3 with \(S=1\) and its selected \(k\).
Exactly the observations outside \([T-kS,T+kS]\) are modified.
If \(m_2\) counts the observations in \((T-kS,T+kS)\), their number is
\(n-m_2\) less the number lying exactly at an endpoint.
Proposal 1 prescribes \(m_2/n\) approximately as \(A(k)\); Proposal 2 determines \(m_2\) through its scale equation;
Proposal 3 determines it through the selected \(k\).
\end{hbtheorem}

This identity is exact for each sample. Huber's further description of
Proposal 1 as asymptotically equivalent to Winsorizing a fixed,
predetermined number of observations, more of them on the side with
heavier contamination, is a description of how the count arises; a
limit theorem comparing Proposal 1 with a fixed-count Winsorized mean is
not supplied here.

\begin{hbtheorem}{thm:p15-w02b-p15-071-prop2-minimax}{proved=annex:p15-w02b-p15-071-prop2-minimax}{Proposal 2 with matched clipping attains the minimax value}
Let \(0<\varepsilon<1\), let \(k_\varepsilon\) be the normalized constant
of \cref{thm:p15-huber-least-favorable-normal}, and let
\(I_*=(1-\varepsilon)A(k_\varepsilon)\). Put
\[
 r_\varepsilon=\varepsilon+(1-\varepsilon)B(k_\varepsilon)/k_\varepsilon^2 .
\]
There is a unique \(k>0\) with \(B(k)/k^2=r_\varepsilon\). Use
Proposal 2 with this \(k\) and \(\beta=B(k)\). Then
\(\varepsilon<r_\varepsilon\), the cutoff \(q_{\max}\) of
\cref{thm:p15-huber-joint-variance-bound} equals \(k_\varepsilon\), and
\(W(q_{\max},\varepsilon)=1/I_*\).

For every symmetric \(H\) with \(\varepsilon H\{0\}<1-r_\varepsilon\),
the limiting truncated risk
\(\lim_{C\to\infty}\lim_{n\to\infty}E_F\min(nT_n^2,C^2)\) exists and is
at most \(1/I_*\). It equals \(1/I_*\) for every symmetric \(H\)
supported on \(\{|x|>k_\varepsilon\}\).
If \(\varepsilon<1-r_\varepsilon\), this covers every symmetric \(H\), and
\[
 \mathcal R_\varepsilon(T)=1/I_*=
 \inf_{T'\ {\rm translation\ equivariant}}\mathcal R_\varepsilon(T'),
\]
with \(\mathcal R_\varepsilon\) of \cref{def:p15-huber-symmetric-location-model}.
For any other \(k\) with \(B(k)/k^2>\varepsilon\), the worst case
\(W(q_{\max},\varepsilon)\) is the worst variance of the known-scale
clipped estimate at cutoff \(q_{\max}\), hence at least \(1/I_*\).
\end{hbtheorem}

\begin{hbexample}{def:p15-w02b-p15-071-prop2-numbers}{Matched clipping constants}
The following values are evaluations of the displayed equations:
\[
\begin{array}{c|rrrrr}
\varepsilon&0.001&0.01&0.05&0.1&0.2\\\hline
k_\varepsilon&2.63288&1.94511&1.39838&1.14017&0.86159\\
r_\varepsilon&0.14289&0.24812&0.40712&0.51632&0.65855\\
\text{matched }k&2.62446&1.90666&1.30000&0.99935&0.67068\\
1-r_\varepsilon&0.85711&0.75188&0.59288&0.48368&0.34145\\
1/I_*&1.00955&1.06524&1.25612&1.48985&2.04553
\end{array}
\]
In each column \(\varepsilon<1-r_\varepsilon\), so the conclusion holds
over the whole symmetric contamination class. For instance, at
\(\varepsilon=0.1\) the matched constant is close to one, and the
\(k=1.0\), \(\varepsilon=0.1\) entry \(1.490\) of the historical
Table IV agrees with \(1/I_*\) to the printed precision.
The matching condition \(\varepsilon<1-r_\varepsilon\) fails at
\(\varepsilon=0.5\), where \(1-r_\varepsilon\simeq0.11388\); a computed
crossing lies near \(\varepsilon\simeq0.26989\). Where it fails, a symmetric contaminant with an
atom at zero of mass at least \((1-r_\varepsilon)/\varepsilon\) is not identified, and minimaxity over the entire
class is not established here.
\end{hbexample}

\begin{hbtheorem}{thm:p15-w02b-p15-071-prop1-population}{proved=annex:p15-w02b-p15-071-prop1-population}{Proposal 1: the population worst case}
Let \(0<\varepsilon<1\), \(k>0\), and for symmetric \(H\) let
\(F=(1-\varepsilon)\Phi+\varepsilon H\). Proposal 1's population
cutoff is a \(q>0\) with \(F((-q,q))=A(k)\); it is unique when it
exists. Define its population variance functional
\[
 V_1(F)=\frac{E_F\min(X^2,q^2)}{A(k)^2},
\]
the regular-domain variance formula of
\cref{thm:p15-huber-joint-regular-limit} evaluated at this cutoff.
\begin{enumerate}
\item If \(A(k)<1-\varepsilon\), let \(q_1\) solve
 \((1-\varepsilon)A(q_1)=A(k)\). Every symmetric \(H\) with such a \(q\)
 satisfies \(q\leq q_1\) and \(V_1(F)\leq W(q_1,\varepsilon)\), with
 equality for every symmetric \(H\) supported on \(\{|x|>q_1\}\).
\item If \(A(k)\geq1-\varepsilon\), then
 \(\sup_HV_1(F)=\infty\) over symmetric \(H\) having such a \(q\).
\item With \(A(k)=(1-\varepsilon)A(k_\varepsilon)\), the worst case is
 \(W(k_\varepsilon,\varepsilon)=1/I_*\), the minimax value.
\end{enumerate}
\end{hbtheorem}

This proves the population form of Huber's minimax statement for
Proposal 1. The matched constants satisfy
\(A(k)=(1-\varepsilon)A(k_\varepsilon)\), giving
\(k\simeq2.59504\), \(1.87173\), \(1.27052\), \(0.97655\), \(0.65708\)
for \(\varepsilon=0.001\), \(0.01\), \(0.05\), \(0.1\), \(0.2\).
The sampling statements need a definite reading of the count equation,
whose left side is an integer divided by \(n\). The following
convention supplies one, and the results after it give Proposal 1's
sampling theory under it.

\begin{hbdefinition}{def:p15-w02b-p15-071-prop1-convention}{Proposal 1 with a closed-window count}
Fix \(k>0\) and integers \(m_n\) with \(1\leq m_n\leq n\) and
\(|m_n-nA(k)|\leq1\), for instance \(m_n=\max(1,\lceil nA(k)\rceil)\).
For a sample \(x_1,\ldots,x_n\) and real \(t\), let
\(D_{(1)}(t)\leq\cdots\leq D_{(n)}(t)\) be the ordered values of
\(|x_i-t|\), and put \(Q_n(t)=D_{(m_n)}(t)\). A pair \((T,S)\) solves
Proposal 1 when
\[
 S>0,\qquad kS=Q_n(T),\qquad \sum_i\psi(k,(x_i-T)/S)=0 .
\]
The condition \(kS=Q_n(T)\) holds exactly when the closed interval
\([T-kS,T+kS]\) contains at least \(m_n\) observations and the open
interval \((T-kS,T+kS)\) contains fewer than \(m_n\). This reads the
source's count of the observations in \([T-kS,T+kS]\)
\cite[p.~96]{Huber1964RobustLocation} as a closed-window order
statistic with target integer \(m_n\).

Let \(E_n\) be the event that no value occurs \(m_n\) or more times in
the sample. On \(E_n\) let \(T_n\) be the smallest \(T\) for which a
solution \((T,S)\) exists, and put \(S_n=Q_n(T_n)/k\). Off \(E_n\) let
\(T_n\) be the midpoint median and \(S_n=1\). The sequence
\(T^{(1)}=(T_n)\) is translation equivariant.

For \(F=(1-\varepsilon)\Phi+\varepsilon H\) with \(0\leq\varepsilon<1\)
and \(H\) symmetric, put
\[
 q_0=\inf\{q>0:F([-q,q])\geq A(k)\},\quad
 \alpha_0=F((-q_0,q_0)),\quad
 p_0=F\{q_0\}=F\{-q_0\},
\]
and \(\sigma_0^2=E_F\min(X^2,q_0^2)\). Call \(F\) regular when
\(p_0=0\). The standing condition is
\[
 \text{(I)}\qquad \max_xF\{x\}<A(k).
\]
Since \(F\{x\}\leq\varepsilon\) for every \(x\), (I) holds for every
symmetric \(H\) when \(\varepsilon<A(k)\). Finally write
\[
 V^*(q,\varepsilon)=\frac{(1-\varepsilon)B(q)+\varepsilon q^2}
                         {(1-\varepsilon)^2A(q)^2},\qquad q>0,
\]
the worst known-scale variance of \cref{thm:p15-huber-normal-variance}
at cutoff \(q\). The bound \(W(q_1,\varepsilon)\) of
\cref{thm:p15-w02b-p15-071-prop1-population} is \(V^*(q_1,\varepsilon)\);
at \(q=q_{\max}\) of \cref{thm:p15-huber-joint-variance-bound} the two
expressions agree because \(B(q_{\max})/q_{\max}^2=(r-\varepsilon)/(1-\varepsilon)\).
\end{hbdefinition}

\begin{hblemma}{thm:p15-w02b-p15-071-empirical-tools}{proved=annex:p15-w02b-p15-071-empirical-tools}{Empirical tools for the count equation}
Let \(X_1,X_2,\ldots\) be i.i.d.\ with law \(F\), let \(F_n\) be the
empirical distribution function of the first \(n\), and
\(D_n=\sup_x|F_n(x)-F(x)|\). Write \(\psi_q(x)=\max(-q,\min(x,q))\).
\begin{enumerate}
\item For every nondecreasing \(g:\mathbb R\to[-1,1]\),
 \(|\int g\,dF_n-\int g\,dF|\leq2D_n\); for every interval \(J\),
 \(|F_n(J)-F(J)|\leq2D_n\).
\item \(\lim_{q\downarrow0}\sup_tF([t-q,t+q])=\max_xF\{x\}\).
\item Let \(F\) be symmetric about zero and \(\ell>0\). Put
 \(R_n(q)=n^{-1/2}\sum_i\{\psi_q(X_i)-\psi_\ell(X_i)\}\). For
 \(\kappa>0\) and \(x>0\),
 \(P(\sup_{|q-\ell|\leq\kappa}|R_n(q)|\geq x)\leq\kappa^2/x^2\).
\item If \(F\) is differentiable at \(x_0\), then for every \(M>0\) the
 supremum of \(\sqrt n|F_n(J)-F(J)|\) over all intervals
 \(J\subseteq[x_0-M/\sqrt n,x_0+M/\sqrt n]\) tends to zero in
 probability.
\item If \(Y_n\Rightarrow N(0,\sigma^2)\) with \(\sigma>0\), and
 \(U_n\to u>0\), \(V_n\to0\) in probability, then
 \((Y_n+V_n)/U_n\Rightarrow N(0,\sigma^2/u^2)\), the quotient being
 given any value where \(U_n\leq0\).
\end{enumerate}
\end{hblemma}

\begin{hbtheorem}{thm:p15-w02b-p15-071-prop1-consistency}{proved=annex:p15-w02b-p15-071-prop1-consistency}{Existence and consistency of every Proposal 1 solution}
Let \(F=(1-\varepsilon)\Phi+\varepsilon H\), \(0\leq\varepsilon<1\),
\(H\) symmetric, \(k>0\), and assume (I). Then \(0<q_0<\infty\),
\(\alpha_0>0\) and
\(\alpha_0\leq A(k)\leq\alpha_0+2p_0\).
Almost surely, for all sufficiently large \(n\), the event \(E_n\)
occurs, Proposal 1 has solutions, and
\[
 \sup\{|T|+|kS-q_0|:(T,S)\text{ solves Proposal 1}\}\longrightarrow0 .
\]
In particular \(T_n\to0\) and \(S_n\to q_0/k\) almost surely.
\end{hbtheorem}

\begin{hbtheorem}{thm:p15-w02b-p15-071-prop1-limit}{proved=annex:p15-w02b-p15-071-prop1-limit}{Limiting risk of Proposal 1, including window atoms}
Under the hypotheses of \cref{thm:p15-w02b-p15-071-prop1-consistency},
let \((T_n',S_n')\) be any measurable choice of solutions on \(E_n\),
for instance \((T_n,S_n)\).
\begin{enumerate}
\item In every case
 \(\lim_{C\to\infty}\limsup_nE\min(nT_n'^2,C^2)\leq\sigma_0^2/\alpha_0^2\).
\item If \(p_0=0\), then \(\alpha_0=A(k)\), \(q_0\) is the population
 cutoff of \cref{thm:p15-w02b-p15-071-prop1-population}, and
 \(\sqrt nT_n'\Rightarrow N(0,V_1(F))\) with
 \(V_1(F)=\sigma_0^2/A(k)^2\).
\item If \(p_0>0\) and \(\alpha_0<A(k)<\alpha_0+p_0\), then almost
 surely, for all large \(n\), every solution has \(kS=q_0-|T|\), and
 \(\sqrt nT_n'\Rightarrow N(0,\sigma_0^2/\alpha_0^2)\).
\item If \(p_0>0\) and \(\alpha_0+p_0<A(k)<\alpha_0+2p_0\), then almost
 surely, for all large \(n\), every solution has \(kS=q_0+|T|\), and
 \(\sqrt nT_n'\Rightarrow N(0,\sigma_0^2/(\alpha_0+2p_0)^2)\).
\item If \(F\) is differentiable at \(q_0\), with derivative
 \(f(q_0)\), then \(p_0=0\), \(f(q_0)>0\) and
 \[
  \sqrt n\{\log S_n'-\log(q_0/k)\}\Longrightarrow
  N\!\left(0,\frac{A(k)(1-A(k))}{4q_0^2f(q_0)^2}\right).
 \]
\end{enumerate}
In (ii)--(iv) the limiting truncated risk
\(\lim_{C\to\infty}\lim_nE\min(nT_n'^2,C^2)\) exists and equals the
variance of the normal limit.
\end{hbtheorem}

Part (iii) is the effect of the count convention at an atom on the
window's edge. When the nearer atom alone carries the count past
\(m_n\), the window shrinks as \(T\) moves, so that it never contains
either atom in its interior. The estimate then uses only the
interior mass \(\alpha_0<A(k)\), and its variance exceeds the regular
formula. In part (iv) the window always contains both atoms and the
variance falls below it. When \(p_0>0\) and \(A(k)\) equals
\(\alpha_0\), \(\alpha_0+p_0\) or \(\alpha_0+2p_0\), the window sits on
an atom boundary at the scale of the sampling fluctuations of the atom
counts relative to \(m_n\). For these three configurations the theorem
gives only the bound (i); their limit law is not determined here.

\begin{hbtheorem}{thm:p15-w02b-p15-071-prop1-minimax}{proved=annex:p15-w02b-p15-071-prop1-minimax}{Proposal 1's asymptotic minimax standing}
Let \(0<\varepsilon<1\) and \(k>0\) with \(A(k)<1-\varepsilon\), and let
\(q_1\) be as in \cref{thm:p15-w02b-p15-071-prop1-population}.
\begin{enumerate}
\item For every symmetric \(H\) satisfying (I), \(q_0\leq q_1\) and
 \(\lim_{C}\limsup_nE\min(nT_n^2,C^2)\leq V^*(q_0,\varepsilon)\).
\item Over regular symmetric \(H\) satisfying (I), the limiting
 truncated risk of \(T^{(1)}\) is at most \(V^*(q_1,\varepsilon)\). It
 equals \(V^*(q_1,\varepsilon)\) for every symmetric \(H\) supported on
 \(\{|x|>q_1\}\) and satisfying (I).
\item Let \(A(k)=(1-\varepsilon)A(k_\varepsilon)\). Then
 \(q_1=k_\varepsilon\), and the bound in (ii) is \(1/I_*\): over regular
 laws satisfying (I), Proposal 1 attains the minimax value. Let
 \(q_*=0\) if \(A(k)\leq\varepsilon/2\), and otherwise let \(q_*\) solve
 \((1-\varepsilon)A(q_*)=A(k)-\varepsilon/2\). For every
 \(q\in(q_*,k_\varepsilon)\) there is a symmetric \(H\) satisfying (I)
 for which configuration (iii) of
 \cref{thm:p15-w02b-p15-071-prop1-limit} holds with \(q_0=q\), and the
 limiting truncated risk is \(V^*(q,\varepsilon)>1/I_*\). Consequently
 \[
  \mathcal R_\varepsilon(T^{(1)})\geq
  \sup_{q_*<q<k_\varepsilon}V^*(q,\varepsilon)>1/I_*,
 \]
 and under this convention Proposal 1 with matched \(k\) is not
 asymptotically minimax over \(\mathcal C_\varepsilon\).
\item If \(\varepsilon<A(k)\), then
 \(\mathcal R_\varepsilon(T^{(1)})\leq\sup_{0<q\leq q_1}V^*(q,\varepsilon)<\infty\).
\end{enumerate}
\end{hbtheorem}

Thus Huber's statement that Proposal 1 is asymptotically minimax when
\(k\) is chosen in relation to the true \(\varepsilon\)
\cite[p.~96]{Huber1964RobustLocation} holds over the regular laws and
fails over the whole symmetric class, the failure coming from
contaminants with atoms at the edge of the counting window. The
corresponding statement for Proposal 2 holds over the whole class when
\(\varepsilon<1-r_\varepsilon\), by
\cref{thm:p15-w02b-p15-071-prop2-minimax}.

\begin{hbexample}{def:p15-w02b-p15-071-prop1-numbers}{The size of the excess at matched clipping}
The following values are evaluations of the equations of
\cref{thm:p15-w02b-p15-071-prop1-minimax}(iii); \(k\) is the matched
Proposal 1 constant and \(V^*(q_*,\varepsilon)\) is a lower bound for
\(\mathcal R_\varepsilon(T^{(1)})\), by continuity of \(V^*\).
\[
\begin{array}{c|rrrrr}
\varepsilon&0.01&0.05&0.1&0.2&0.3\\\hline
k&1.87173&1.27052&0.97655&0.65708&0.45992\\
A(k)&0.93876&0.79610&0.67121&0.48887&0.35443\\
q_*&1.90475&1.31558&1.01570&0.65276&0.37460\\
V^*(q_*,\varepsilon)&1.06533&1.25725&1.49405&2.06690&2.89589\\
1/I_*&1.06524&1.25612&1.48985&2.04553&2.82145
\end{array}
\]
In each column \(\varepsilon<A(k)\), so (I) holds for every symmetric
\(H\); a computed crossing of \(\varepsilon=A(k)\) lies near
\(\varepsilon\simeq0.32545\). For \(\varepsilon=0.1\) and \(q=1.05\),
take \(H=h(\delta_{1.05}+\delta_{-1.05})+(1-2h)U\) with \(h=0.45\) and
\(U\) uniform on \([-3.05,-2.05]\cup[2.05,3.05]\). Then
\(\alpha_0\simeq0.63565<A(k)<\alpha_0+p_0\simeq0.68065\), and the
limiting truncated risk is \(V^*(1.05,0.1)\simeq1.49204\), against
\(\sigma_0^2/A(k)^2\simeq1.338\) for the regular formula. The excess
over \(1/I_*\) is small in these columns; its sign, not its size, is
what the theorem establishes.

At a \(k\) not adapted to \(\varepsilon\), the regular laws of
\cref{thm:p15-w02b-p15-071-prop1-population}(ii), with an atom at zero
and far symmetric atoms, satisfy (I) when \(\varepsilon<A(k)\). By
\cref{thm:p15-w02b-p15-071-prop1-limit}(ii) their limiting truncated
risks equal \(V_1(F)\), so the supremum is infinite at the sampling
level as well. At \(k=2\) and \(\varepsilon=0.05\), \(A(2)\simeq0.9545\)
exceeds \(1-\varepsilon\), and Proposal 1's worst limiting risk is
unbounded, whereas the corresponding Table IV entry for Proposal 2 is
\(1.353\).
\end{hbexample}

\begin{hbtheorem}{thm:p15-w02b-p15-071-scale-compare}{proved=annex:p15-w02b-p15-071-scale-compare}{The scale accuracy comparison at the normal law}
Let the observations be standard normal and \(k>0\), and put
\[
 C(k)=A(k)-2k\phi(k),\qquad
 M_4(k)=3A(k)-2(k^3+3k)\phi(k)+k^4(1-A(k)).
\]
\begin{enumerate}
\item For Proposal 1 under \cref{def:p15-w02b-p15-071-prop1-convention},
 \(\sqrt n\log S_n\Rightarrow N(0,v_1(k))\), with
 \(v_1(k)=A(k)(1-A(k))/\{4k^2\phi(k)^2\}\).
\item For Proposal 2 with \(\beta=B(k)\),
 \(\sqrt n\log S_n\Rightarrow N(0,v_2(k))\), with
 \(v_2(k)=\{M_4(k)-B(k)^2\}/\{4C(k)^2\}\).
\item \(v_1(1)<0.93<1.02<v_2(1)\) and \(v_2(2)<0.57<0.92<v_1(2)\).
\end{enumerate}
\end{hbtheorem}

Huber states that the \(S\) of Proposal 2 will be more accurate in
general, \(\log S\) having the smaller asymptotic variance
\cite[p.~96]{Huber1964RobustLocation}. Part (iii) shows the claim
holds at the normal law for \(k=2\) and fails there for \(k=1\); since
\(\Phi\) belongs to every \(\mathcal C_\varepsilon\), the claim does not
hold in general. Evaluations of the two formulas give
\[
\begin{array}{c|rrrr}
k&0.5&1&1.5&2\\\hline
v_1(k)&1.90636&0.92496&0.76677&0.93117\\
v_2(k)&2.32516&1.03350&0.68945&0.56486
\end{array}
\]
and a computed crossing near \(k\simeq1.2974\), found by a grid search on
\(0.01\leq k\leq6\); that this is the only crossing is not proved here.
The matched Proposal 2 constant at \(\varepsilon=0.1\), about \(0.99935\),
lies below that crossing. A comparison under contaminated laws other
than \(\Phi\) is not supplied; part (v) of
\cref{thm:p15-w02b-p15-071-prop1-limit} and
\cref{thm:p15-huber-joint-law-and-variance} give the two variances at
laws without atoms at the respective window edges, but they are not
compared here.

Huber's remaining assessments have the following standing. He presents
all three procedures as heuristic. He conjectures that Proposal 3 is
asymptotically minimax independently of \(\sigma\) and \(\varepsilon\),
and that it has the best asymptotic properties of the three; neither is
proved here. His judgement that Proposal 2 is better than Proposal 1
for large \(n\) and less sensitive to a \(k\) not adapted to
\(\varepsilon\) is supported in the two exact senses given above: at
matched constants Proposal 2 attains the minimax value over the whole
class when \(\varepsilon<1-r_\varepsilon\), while Proposal 1 exceeds it at
window atoms; and at a wrong \(k\), Proposal 1's worst limiting risk can
be infinite where Table IV is finite. A comparison of their small-sample
behavior is not established here. His stated preference for Proposal 2
rests on its fit with iterative least-squares computation; it is a
practical recommendation, not a theorem
\cite[pp.~96--97]{Huber1964RobustLocation}.

\section{The pair-average median under every symmetric contaminant}
\label{sec:p15-w02b-p15-071-pair-median}

\begin{hbtheorem}{thm:p15-w02b-p15-071-pair-median-risk}{proved=annex:p15-w02b-p15-071-pair-median-risk}{The source bound for arbitrary symmetric contamination}
Let \(0\leq\varepsilon<1\), let \(H\) be any probability symmetric about
zero, put \(F=(1-\varepsilon)\Phi+\varepsilon H\), and let \(T_n\) be the
midpoint median of the pair averages \((X_i+X_j)/2\), \(i<j\), of
\(n\geq2\) i.i.d.\ observations from \(F\). Then
\[
 \lim_{C\to\infty}\limsup_{n\to\infty}E_F\min(nT_n^2,C^2)
 \leq\frac{\pi}{3(1-\varepsilon)^4}.
\]
More precisely:
\begin{enumerate}
\item If \(\varepsilon>0\) and \(H\) has an atom, then
 \(P(T_n\neq0)\to0\) and the left side is zero.
\item If \(\varepsilon=0\) or \(H\) has no atom, put
 \[
  \kappa=\frac{(1-\varepsilon)^2}{\sqrt\pi}
          +4\varepsilon(1-\varepsilon)\int\phi\,dH .
 \]
 For every \(z>0\),
 \(\limsup_nP(\sqrt nT_n>z)=\limsup_nP(\sqrt nT_n<-z)\leq
 \Phi(-\sqrt3\,\kappa z)\), and the left side is at most
 \(1/(3\kappa^2)\).
\end{enumerate}
For \(0<\varepsilon<1\), the bound is strict for every symmetric \(H\).
The supremum over all symmetric \(H\) equals the bound.
If \(\sqrt nT_n\) converges in law, the limit's second moment is the
left side and obeys the same bound.
\end{hbtheorem}

This is the source's sharp bound for arbitrary symmetric \(H\)
\cite[pp.~100--101]{Huber1964RobustLocation}, stated for the limiting
truncated risk of \cref{def:p15-huber-symmetric-location-model}.
\Cref{thm:p15-huber-pair-median} gives the normal limit and the formula
\(1/\{12(\int f^2)^2\}\) for square-integrable densities. For a singular
or non-square-integrable contaminant the distribution function of
\(X_1+X_2\) need not be differentiable at zero, and no limit law is
asserted; the risk bound holds regardless. An atom in a symmetric
contaminant produces an atom of \(X_1+X_2\) at zero, which pins the
estimate to zero with probability tending to one.
